\documentclass[11pt,a4paper]{article}
\usepackage{brainteaser}

\begin{document}
\brainteasertitle{5}{Les 25 chevaux}

\begin{statementbox}
On dispose de $25$ chevaux et de $5$ pistes. Une course peut donc opposer au
maximum $5$ chevaux. Il n'y a pas de chronomètre: après chaque course, on
connaît seulement l'ordre d'arrivée. Chaque cheval court toujours à la même
vitesse et il n'y a pas d'égalité.

Quel est le nombre minimal de courses nécessaire pour identifier avec
certitude les trois chevaux les plus rapides? Donner une stratégie et prouver
qu'on ne peut pas faire mieux.
\end{statementbox}

\begin{hintsbox}
\begin{enumerate}
  \item Former cinq groupes de cinq et ordonner chaque groupe.
  \item Quels chevaux faut-il faire courir lors de la sixième course?
  \item Si les vainqueurs sont classés $A_1>B_1>C_1>D_1>E_1$, quels chevaux
  peuvent encore appartenir au top 3?
  \item Pourquoi $B_2$ reste-t-il candidat alors que $C_2$ est éliminé?
  \item Pour la borne inférieure, montrer que cinq courses sont nécessaires
  pour faire courir tous les chevaux, puis qu'une sixième doit comparer les
  cinq vainqueurs.
\end{enumerate}
\end{hintsbox}

\correction

On organise cinq premières courses:
\[
A_1>A_2>A_3>A_4>A_5,
\qquad\ldots\qquad,
E_1>E_2>E_3>E_4>E_5.
\]

La sixième course oppose les cinq vainqueurs. Supposons qu'elle donne
\[
A_1>B_1>C_1>D_1>E_1.
\]
Alors $A_1$ est le plus rapide. Les seuls candidats possibles aux places 2 et
3 sont
\[
A_2,\ A_3,\ B_1,\ B_2,\ C_1.
\]

En effet, $A_4$ a déjà trois chevaux devant lui dans son groupe; $B_3$ est
battu par $B_1,B_2,A_1$; $C_2$ est battu par $C_1,B_1,A_1$; et tout cheval
des groupes $D$ et $E$ a au moins trois supérieurs connus.

Une septième course entre les cinq candidats donne directement les deuxième
et troisième chevaux.

Pour la borne inférieure, tout cheval doit avoir couru au moins une fois,
sinon un cheval non testé pourrait être le plus rapide. Les cinq premières
courses doivent donc être cinq groupes disjoints de cinq. Leurs cinq vainqueurs
sont encore invaincus et incomparables: une sixième course doit les opposer
pour identifier le plus rapide. Elle ne compare toutefois pas tous les
candidats restants aux places 2 et 3, par exemple $A_2$ et $B_1$. Une course
supplémentaire est indispensable.

\begin{resultbox}
\[
\boxed{7\text{ courses sont nécessaires et suffisantes}.}
\]
\end{resultbox}

\end{document}
